\section{案例 4：睡眠作息崩坏}

\subsection*{情境}
一个人想要更早入睡，却在夜间持续消费高刺激内容，结果睡前时间一再后移。

\subsection*{框架映射}
\begin{itemize}[nosep]
  \item \textbf{当前状态：} 高刺激的深夜刷屏或媒体消费。
  \item \textbf{目标状态：} 一个被设计好的睡前机制，并且它能可靠地过渡进睡眠。
  \item \textbf{主导变量：} 高 $\Reward$、延续下去的阻力很低、目标状态规定得很弱，以及让亮度与刺激持续激活的环境。
  \item \textbf{错过的决策窗口：} 在奖赏流完全俘获整晚之前，最早出现的那一个关机线索。
\end{itemize}

\subsection*{机制解释}
这个案例与研究主干之间具有最强的直接连接。关于睡眠转变的文献给出了一个真实例子，说明有模式的状态变化可以用动力学语言来描述。即便如此，手稿仍需保持纪律：被引用的睡眠研究支持的是“转变结构与控制参数很重要”，而不是一整套关于个人睡前纪律的完整行为理论。即便如此，这个案例仍然强烈受益于机制更替的思路：如果亮屏、新奇和无限延展的内容让系统留在错误的盆地里，那么“你就停下来吧”并不是一条真实的控制策略。

\subsection*{证据车道纪律}
在当前来源集中，这个案例拥有最强的\textbf{强证据}锚点，因为被索引的 Nature Neuroscience 睡眠论文直接把入睡处理成一种有模式的转变动力学。把这一结果迁移到实际睡前习惯设计时，仍然有一部分属于\textbf{中等}层级，因为手稿依然是在从生物层面的机制转变，跨到日常行为控制设计。这种跨越是可以辩护的，但它应被表述为解释，而不是已经完成的证明。

\subsection*{操作性修复}
\begin{enumerate}[nosep]
  \item 用一个积极的睡前状态，替代“停止消费”的消极规则。
  \item 明确对象集合：调暗灯光、纸质书、远离床位的充电器，以及一条固定的关机顺序。
  \item 把第一个关机线索视为关键决策窗口，而不是最后那个已经疲惫透支的分钟。
  \item 通过让时间、光线和设备摆放在多个夜晚保持稳定，来减少参数漂移。
\end{enumerate}

\subsection*{迁移教训}
目标状态越是被定义成“坏状态的缺失”，它的进入动力学就越弱。有效控制通常需要一个替代性机制，而不只是一个禁止令。